\section{Main results}\label{sec:main}

We present our framework in three parts.  First, we formalize the saturation
model for CoT verification and analyze its convergence properties.  Second, we
prove bounds connecting proof complexity to verification accuracy and required
compute.  Third, we describe computational experiments validating the
theoretical predictions.

\subsection{The saturation model}\label{sec:saturation}

We model the accuracy of a CoT verifier as an exponential saturation process.

\begin{definition}[CoT verification process]\label{def:cot-process}
For an unsatisfiable CNF formula~$F$, the \emph{CoT verification accuracy} at
reasoning step~$t$ is
\begin{equation}\label{eq:saturation}
p_F(t) = p_{\max}(F) - \bigl(p_{\max}(F) - p_0(F)\bigr) \cdot e^{-\lambda(F) \cdot t},
\end{equation}
where:
\begin{enumerate}[label=(\roman*)]
\item $p_0(F) \in (0.5, 1)$ is the \emph{initial accuracy} (accuracy from a
  single-step assessment);
\item $p_{\max}(F) \in (p_0(F), 1)$ is the \emph{asymptotic accuracy ceiling};
\item $\lambda(F) > 0$ is the \emph{convergence rate}.
\end{enumerate}
\end{definition}

The saturation form~\eqref{eq:saturation} is motivated by the empirical
observation that verification accuracy follows diminishing-returns curves as
a function of test-time compute~\cite{snell2024scaling, shi2025heimdall}.

\begin{definition}[Convergence step]\label{def:convergence-step}
For $\varepsilon > 0$, the \emph{$\varepsilon$-convergence step} of~$F$ is
\[
\tau_F(\varepsilon) = \min\bigl\{t \in \mathbb{N} :
  p_{\max}(F) - p_F(t) < \varepsilon\bigr\}.
\]
\end{definition}

\begin{theorem}[Convergence rate]\label{thm:convergence}
Let $F$ be an unsatisfiable CNF formula and let $p_F(t)$ be defined
by~\eqref{eq:saturation} with parameters $p_0(F)$, $p_{\max}(F)$, and
$\lambda(F) > 0$.  Then:
\begin{enumerate}[label=(\alph*)]
\item $p_F(t)$ is strictly increasing and $\lim_{t \to \infty} p_F(t)
  = p_{\max}(F)$.
\item The $\varepsilon$-convergence step satisfies
\begin{equation}\label{eq:convergence-step}
\tau_F(\varepsilon) = \left\lceil
  \frac{1}{\lambda(F)} \ln\!\left(
  \frac{p_{\max}(F) - p_0(F)}{\varepsilon}\right)\right\rceil.
\end{equation}
\item If $\lambda(F) = \lambda_0 / (1 + C(F))$ for some composite complexity
  measure $C(F) \ge 0$ and constant $\lambda_0 > 0$, then $\tau_F(\varepsilon)$
  is increasing in $C(F)$.
\end{enumerate}
\end{theorem}

\begin{proof}
\textbf{Part~(a).}  We have
$p_F'(t) = (p_{\max}(F) - p_0(F)) \cdot \lambda(F) \cdot e^{-\lambda(F) t}$.
Since $p_{\max}(F) > p_0(F)$ and $\lambda(F) > 0$, the derivative is strictly
positive for all $t \ge 0$.  The limit follows because
$e^{-\lambda(F) t} \to 0$.

\textbf{Part~(b).}  The gap at step~$t$ is
\[
p_{\max}(F) - p_F(t) = (p_{\max}(F) - p_0(F)) \cdot e^{-\lambda(F) t}.
\]
Setting this less than~$\varepsilon$ and solving for~$t$ gives
$t > \frac{1}{\lambda(F)} \ln\!\bigl(\frac{p_{\max}(F) - p_0(F)}{\varepsilon}
\bigr)$.  Taking the ceiling yields~\eqref{eq:convergence-step}.

\textbf{Part~(c).}  Under the stated form,
$\tau_F(\varepsilon) = \lceil \frac{1 + C(F)}{\lambda_0}
\ln\!\bigl(\frac{p_{\max}(F) - p_0(F)}{\varepsilon}\bigr) \rceil$.
Since $1 + C(F)$ is increasing in $C(F)$ and the logarithmic factor is
positive, $\tau_F(\varepsilon)$ is increasing in~$C(F)$.
\end{proof}

We now specify how the convergence rate depends on proof structure.

\begin{definition}[Composite complexity measure]\label{def:composite}
For an unsatisfiable CNF formula~$F$, the \emph{composite complexity measure}
is
\begin{equation}\label{eq:composite}
C(F) = \log S(F) + \log w(F \vdash \bot) + \log d(F),
\end{equation}
where $S(F)$, $w(F \vdash \bot)$, and $d(F)$ are the minimum proof length,
width, and depth, respectively.
\end{definition}

\begin{remark}\label{rem:composite}
The choice of a sum of logarithms corresponds to taking the logarithm of the
product $S(F) \cdot w(F \vdash \bot) \cdot d(F)$.  This multiplicative
combination captures the intuition that proof verification difficulty grows
with each dimension of proof complexity.  The logarithmic scaling ensures that
the composite measure is comparable across formula families with vastly
different parameter ranges.
\end{remark}

\subsection{Error accumulation and accuracy bounds}\label{sec:error}

We now prove that proof length governs the asymptotic accuracy ceiling through
an error accumulation mechanism.

\begin{theorem}[Error accumulation bound]\label{thm:error-accumulation}
Let $F$ be an unsatisfiable CNF formula with minimum proof length $S(F)$.
Suppose that a CoT verifier checks each proof step independently, accepting
step~$i$ correctly with probability at least $1 - \varepsilon_0$ for some
$\varepsilon_0 \in (0, 1/2)$.  Then the probability that the verifier
correctly identifies the proof as valid satisfies
\begin{equation}\label{eq:error-bound}
p_{\max}(F) \le (1 - \varepsilon_0)^{S(F)} \le e^{-\varepsilon_0 \cdot S(F)}.
\end{equation}
\end{theorem}

\begin{proof}
Under the step-independence assumption, the verifier accepts the proof as valid
if and only if it accepts every step.  Since each step is accepted correctly
with probability at most $1 - \varepsilon_0$, the probability that all $S(F)$
steps are accepted correctly is at most $(1 - \varepsilon_0)^{S(F)}$.  The
second inequality follows from $1 - x \le e^{-x}$ for all $x \ge 0$.
\end{proof}

\begin{corollary}\label{cor:accuracy-length}
Under the hypotheses of Theorem~\ref{thm:error-accumulation}, if
$S(F) \ge \ln(1/\delta) / \varepsilon_0$, then $p_{\max}(F) \le \delta$.
\end{corollary}

\begin{proof}
From~\eqref{eq:error-bound},
$p_{\max}(F) \le e^{-\varepsilon_0 \cdot S(F)} \le e^{-\ln(1/\delta)} =
\delta$.
\end{proof}

\begin{remark}\label{rem:step-independence}
The step-independence assumption is a simplification.  In practice, errors at
one step may be correlated with errors at subsequent steps---for example,
because the CoT trace conditions on prior reasoning.  Ton et
al.~\cite{ton2025understanding} show that once a model encounters an
unidentifiable sub-task, all subsequent steps contribute zero information,
which is a stronger form of error propagation.  Theorem~\ref{thm:error-accumulation}
therefore provides an \emph{optimistic} bound; the true accuracy degradation
with proof length may be worse.
\end{remark}

\begin{example}\label{ex:php-accuracy}
Consider the pigeonhole formula $\operatorname{PHP}^{n+1}_n$.  By
Haken~\cite{haken1985intractability}, $S(\operatorname{PHP}^{n+1}_n) =
2^{\Omega(n)}$.  Even with a highly accurate per-step verifier having
$\varepsilon_0 = 0.001$, Theorem~\ref{thm:error-accumulation} gives
$p_{\max} \le e^{-0.001 \cdot 2^{\Omega(n)}}$, which drops below any
constant for sufficiently large~$n$.  This shows that pigeonhole formulas
present a fundamental barrier for step-by-step verification.
\end{example}

\subsection{Self-consistency amplification}\label{sec:amplification}

We analyze how majority voting over independent chains can boost verification
accuracy, and connect the required number of chains to proof complexity.

\begin{theorem}[Self-consistency amplification]\label{thm:self-consistency}
Let $Y_1, \ldots, Y_N$ be independent Bernoulli random variables with common
parameter $p > 1/2$.  The majority vote
$\hat{Y}_N = \mathbf{1}[\sum_{i=1}^N Y_i > N/2]$ satisfies:
\begin{enumerate}[label=(\alph*)]
\item $\Pr[\hat{Y}_N = 1] \ge 1 - \exp\bigl(-2N(p - \tfrac{1}{2})^2\bigr)$.
\item To achieve $\Pr[\hat{Y}_N = 1] \ge 1 - \delta$, it suffices to take
\begin{equation}\label{eq:chains-needed}
N \ge \frac{\ln(1/\delta)}{2(p - \tfrac{1}{2})^2}.
\end{equation}
\end{enumerate}
\end{theorem}

\begin{proof}
\textbf{Part~(a).}  Let $\bar{Y} = \frac{1}{N}\sum_{i=1}^N Y_i$.  The
majority vote errs when $\bar{Y} \le 1/2$.  By Hoeffding's
inequality~\cite{hoeffding1963probability}, for any $t > 0$,
\[
\Pr\!\left[\bar{Y} - p \le -t\right]
  \le \exp(-2Nt^2).
\]
Setting $t = p - 1/2 > 0$, we obtain
\[
\Pr\!\left[\bar{Y} \le \tfrac{1}{2}\right]
  \le \Pr\!\left[\bar{Y} - p \le -(p - \tfrac{1}{2})\right]
  \le \exp\!\left(-2N(p - \tfrac{1}{2})^2\right).
\]
Therefore $\Pr[\hat{Y}_N = 1] \ge 1 -
\exp\bigl(-2N(p - \tfrac{1}{2})^2\bigr)$.

\textbf{Part~(b).}  Setting $\exp(-2N(p - 1/2)^2) \le \delta$ and solving
for~$N$ gives~\eqref{eq:chains-needed}.
\end{proof}

\begin{corollary}[Complexity--compute tradeoff]\label{cor:compute-tradeoff}
Suppose the base verification accuracy satisfies $p(F) > 1/2$ and depends on
the composite complexity $C(F)$ through a decreasing function
$p(F) = g(C(F))$ with $g(C) \to 1/2$ as $C \to \infty$.  Then the number of
self-consistency chains needed to reach accuracy $1 - \delta$ grows as
$\Omega(1/(g(C(F)) - 1/2)^2)$, which diverges as $C(F) \to \infty$.
\end{corollary}

\begin{proof}
By Theorem~\ref{thm:self-consistency}(b), $N \ge \ln(1/\delta) /
(2(p(F) - 1/2)^2)$.  Since $p(F) = g(C(F)) \to 1/2$ as $C(F) \to \infty$,
the denominator vanishes, so $N \to \infty$.
\end{proof}

\begin{remark}\label{rem:hoeffding-tightness}
The Hoeffding bound is known to be loose for Bernoulli random variables; the
exact majority probability can be computed via the binomial CDF.  However, the
Hoeffding bound correctly captures the \emph{scaling}: $N = \Theta(1/(p -
1/2)^2)$ is tight up to constant factors.  Our computational experiments
confirm that the ordering of chain requirements across complexity levels is
perfectly preserved ($\rho = 1.0$), even though absolute values differ.
\end{remark}

\begin{example}\label{ex:chains-needed}
Consider two formulas: a Tseitin formula on a path graph with base accuracy
$p = 0.90$ and a pigeonhole formula with $p = 0.60$.  To reach $95\%$ majority
vote accuracy ($\delta = 0.05$), Theorem~\ref{thm:self-consistency} requires:
\begin{align*}
N_{\text{Tseitin}} &\ge \frac{\ln 20}{2(0.40)^2} \approx 9.4
  \implies N = 10, \\
N_{\text{PHP}} &\ge \frac{\ln 20}{2(0.10)^2} \approx 149.8
  \implies N = 150.
\end{align*}
The pigeonhole formula requires roughly $15\times$ more chains, reflecting its
greater proof complexity.
\end{example}

\subsection{Phase transition in verification difficulty}\label{sec:phase}

We provide a formal criterion for the existence of a phase transition in
verification accuracy as a function of proof complexity.

\begin{definition}[Piecewise linear model]\label{def:piecewise}
A \emph{piecewise linear model} for verification accuracy as a function of
composite complexity $C$ is
\begin{equation}\label{eq:piecewise}
p_{\max}(C) = \begin{cases}
  a_1 + b_1 \cdot C & \text{if } C \le C^*, \\
  a_2 + b_2 \cdot C & \text{if } C > C^*,
\end{cases}
\end{equation}
where $C^*$ is the \emph{critical threshold}, $b_1, b_2 < 0$ are degradation
rates, and continuity at $C^*$ imposes $a_1 + b_1 C^* = a_2 + b_2 C^*$.
\end{definition}

\begin{proposition}[Phase transition criterion]\label{prop:phase-transition}
If the piecewise linear model~\eqref{eq:piecewise} provides a statistically
significant improvement over a single linear model $p_{\max}(C) = a + bC$
(as measured by an $F$-test with $p < 0.05$), and $\abs{b_2} > \abs{b_1}$,
then we say a \emph{phase transition in verification difficulty} occurs at
$C^*$.
\end{proposition}

\begin{remark}\label{rem:phase-analogy}
This notion of phase transition is analogous to the well-known phase
transition in random $k$-SAT, where the probability of satisfiability changes
sharply at a critical clause-to-variable ratio~\cite{mezard2002analytic}.
In our setting, the ``order parameter'' is verification accuracy rather than
satisfiability, and the control parameter is proof complexity rather than
clause density.
\end{remark}

\subsection{Computational verification}\label{sec:computation}

We validate the theoretical predictions of
Sections~\ref{sec:saturation}--\ref{sec:phase} on a dataset of~129 synthetic
proof instances.

\subsubsection{Experimental setup}\label{sec:setup}

We generate proof instances from four formula families with known complexity
parameters:

\begin{table}[ht]
\centering
\caption{Summary of the 129 proof instances used for computational validation.
The complexity range gives the minimum and maximum of the composite measure
$C(F) = \log S(F) + \log w(F \vdash \bot) + \log d(F)$ within each
family.}\label{tab:instances}
\begin{tabular}{@{}lccc@{}}
\toprule
\textbf{Family} & \textbf{Count} & \textbf{Graph types / parameters}
  & \textbf{$C(F)$ range} \\
\midrule
Tseitin & 27 & Path, cycle, grid, expander & $[1.6, 14.2]$ \\
Pigeonhole & 9 & $\operatorname{PHP}^{n+1}_n$, $n = 3, \ldots, 11$
  & $[5.3, 19.8]$ \\
Pebbling & 18 & Chain, tree, pyramid DAGs & $[3.2, 18.5]$ \\
Random 3-CNF & 75 & Ratios $\alpha \in [2.0, 5.0]$ & $[1.0, 12.7]$ \\
\bottomrule
\end{tabular}
\end{table}

For each instance~$F$, we compute exact complexity measures for the structured
families and estimates for random $3$-CNF\@.  We then simulate the CoT
verification process~\eqref{eq:saturation} with 200--500 trials per instance,
using a convergence rate model
\begin{equation}\label{eq:lambda-model}
\lambda(F) = \frac{\lambda_0}{1 + \sum_{i} w_i \cdot C_i(F)},
\end{equation}
where $C_i(F)$ are individual log-normalized complexity features and
$w_i$ are weights calibrated to published empirical data: $w_{\text{width}} =
0.35$, $w_{\text{space}} = 0.25$, $w_{\text{length}} = 0.25$,
$w_{\text{depth}} = 0.15$, and $\lambda_0 = 2.0$.  The calibration targets
are: (i)~single-pass accuracy of $94.5\%$ and $64$-vote accuracy of $97.5\%$
from Heimdall~\cite{shi2025heimdall}; (ii)~difficulty-dependent saturation
curves from Snell et al.~\cite{snell2024scaling}.

\subsubsection{Results}

\paragraph{Complexity--convergence correlation.}
Table~\ref{tab:correlations} shows Spearman rank correlations between
complexity features and convergence metrics.

\begin{table}[ht]
\centering
\caption{Spearman rank correlations between complexity features and
convergence rate~$\lambda$.  All correlations are significant at the
$p < 10^{-10}$ level after Bonferroni correction.}\label{tab:correlations}
\begin{tabular}{@{}lcc@{}}
\toprule
\textbf{Feature} & \textbf{Spearman $\rho$} & \textbf{$p$-value} \\
\midrule
$C(F)$ (composite) & $-0.991$ & $< 10^{-113}$ \\
$\log S(F)$ (length) & $-0.840$ & $< 10^{-35}$ \\
$\log \operatorname{Sp}(F)$ (space) & $-0.782$ & $< 10^{-27}$ \\
$\log d(F)$ (depth) & $-0.678$ & $< 10^{-18}$ \\
$\log w(F \vdash \bot)$ (width) & $-0.634$ & $< 10^{-15}$ \\
\bottomrule
\end{tabular}
\end{table}

The composite measure~$C(F)$ achieves the strongest correlation
($\rho = -0.991$), consistent with Theorem~\ref{thm:convergence}(c).  Among
individual measures, proof length is the strongest single predictor, followed
by space.  Multiple linear regression on all complexity features yields
$R^2 = 0.93$.

\paragraph{Accuracy--length relationship.}
The correlation between asymptotic accuracy $p_{\max}(F)$ and proof length is
$\rho = -0.766$ ($p < 10^{-25}$), consistent with
Theorem~\ref{thm:error-accumulation}.  Accuracy ranges from~$0.972$ for simple
Tseitin formulas on path graphs to~$0.880$ for large pigeonhole instances.

\paragraph{Phase transition.}
Fitting the piecewise linear model~\eqref{eq:piecewise} yields a critical
threshold $C^* \approx 12.95$.  The $F$-test comparing the piecewise model to
a single linear fit gives $F = 7.24$ with $p = 0.008$.  Below~$C^*$, accuracy
degrades at rate $b_1 = -0.0054$ per unit complexity; above~$C^*$, the rate
steepens to $b_2 = -0.0069$, a factor of~$1.28$.  This confirms the phase
transition criterion of Proposition~\ref{prop:phase-transition}.

\paragraph{Self-consistency chains.}
Table~\ref{tab:chains} compares empirically required chains with the Hoeffding
prediction of Theorem~\ref{thm:self-consistency}.

\begin{table}[ht]
\centering
\caption{Number of independent chains needed to reach $95\%$ majority vote
accuracy.  The Hoeffding bound is loose but preserves the ordering perfectly
(Spearman $\rho = 1.0$).}\label{tab:chains}
\begin{tabular}{@{}ccc@{}}
\toprule
\textbf{Base accuracy $p$} & \textbf{Chains (empirical)}
  & \textbf{Chains (Hoeffding)} \\
\midrule
$0.60$ & $67$ & $299$ \\
$0.65$ & $29$ & $133$ \\
$0.70$ & $17$ & $75$ \\
$0.75$ & $9$ & $48$ \\
$0.80$ & $7$ & $33$ \\
$0.85$ & $5$ & $24$ \\
$0.90$ & $3$ & $15$ \\
\bottomrule
\end{tabular}
\end{table}

\paragraph{Per-family comparison.}
Table~\ref{tab:families} summarizes convergence metrics by formula family.

\begin{table}[ht]
\centering
\caption{Convergence metrics by formula family ($\text{mean} \pm
\text{std}$).}\label{tab:families}
\begin{tabular}{@{}lccc@{}}
\toprule
\textbf{Family} & \textbf{Conv.\ step} & \textbf{Conv.\ rate $\lambda$}
  & \textbf{Final accuracy} \\
\midrule
Tseitin & $4.9 \pm 4.5$ & $0.664 \pm 0.259$ & $0.940 \pm 0.023$ \\
Pebbling & $6.2 \pm 4.8$ & $0.612 \pm 0.272$ & $0.941 \pm 0.021$ \\
Random 3-CNF & $4.8 \pm 4.5$ & $0.570 \pm 0.087$ & $0.928 \pm 0.016$ \\
Pigeonhole & $6.2 \pm 4.6$ & $0.514 \pm 0.176$ & $0.923 \pm 0.028$ \\
\bottomrule
\end{tabular}
\end{table}

Pigeonhole formulas exhibit the lowest convergence rate and accuracy,
consistent with their exponential proof complexity lower
bounds~\cite{haken1985intractability}.  Tseitin formulas on simple graphs are
the easiest, reflecting their polynomial-size proofs on low-expansion graphs.
